%% Q1
The most important quality of a good teacher is the ability to \textbf{motivate students to learn}. Motivation drives attention, effort and lifelong learning; mastery of subject and methods serve this end.

%% Q2
During teaching, a teacher should \textbf{allow students to raise questions}. Questions show engagement, clarify doubts and make learning two-way.

%% Q3
Per the key, the ultimate goal of education is to \textbf{develop problem-solving skills} so that learners can face real-life situations independently.

%% Q4
The main functions of teachers in higher education are \textbf{teaching, research and extension} (service to the community), the three pillars of universities.

%% Q5
An honest teacher \textbf{accepts the lack of knowledge} and promises to come prepared next time, modelling integrity and a learning attitude.

%% Q6
Scientific research aims at theory building, problem solving and creating new knowledge. \textbf{Logical reasoning} is a \emph{means} (a tool) of research, not its aim.

%% Q7
Studying the same subjects repeatedly over several years is a \textbf{longitudinal study}. A cross-sectional study observes different groups at one point in time.

%% Q8
The most important factor is \textbf{researchability}: whether the problem can be studied empirically with available data, tools, time and resources.

%% Q9
\textbf{Inductive} reasoning moves from \textbf{particular to general} (specific observations $\rightarrow$ general conclusion). Deduction moves from general to particular.

%% Q10
Case study, ethnography and content analysis are qualitative. A \textbf{pretest--posttest experiment} is a quantitative experimental design.

%% Q11 ctx
\textbf{Direction (Q.~No.~11 to 15):} Read the following text and answer questions 11 to 15.\par\smallskip
Sixty per cent of rural India lives in utter poverty with no electricity. Rural India uses 180 million tones of fuel every year for cooking, creating indoor pollution and health hazards. The WHO estimates that 1.5 million deaths per year are caused by pollution from chulha smoke. Making available locally grown clean and renewable cooking and lighting fuel can improve the quality of rural life. Liquid fuels are far superior to solid fuels for cooking because of their clean burning and higher energy. Ethanol is one of the best as it is an excellent substitute for kerosene and burns better with no unpleasant smell. Its burning power is almost as clean as that of LPG.

Ethanol can be produced from any sugary material and is presently produced world over from sugarcane and corn. However, as the ethanol economy grows, there is need to produce it from a crop which uses much less water than sugarcane and also produces food. Sorghum (Jowar) is one such crop whose stem is sweet like sugarcane. Its earhead produces grains which can be used for making bread. Its sweet stem has nearly the same amount of sugar and hence the juice can be fermented and used for ethanol production. The left-over stem after juice extraction, together with leaves is excellent fodder for animals. So from the same price of land one can get food, fuel and fodder. Besides, sweet sorghum (Jowar) uses nearly fifty percent less water than sugarcane to produce the same amount of sugar. It is a four month crop so farmers can grow two crops per year from the same piece of land. Also the energy output ratio from sweet sorghum is very positive.

%% Q11 q
The purpose of the paragraph is to bring out:

%% Q11
The paragraph identifies a problem (polluting solid cooking fuels) and proposes a solution: clean, locally grown ethanol from sweet sorghum. Its purpose is to bring out \textbf{remedies for pollutant energy}.

%% Q12
The passage argues for \textbf{alternate fuels} (ethanol from sweet sorghum) in place of polluting solid fuels.

%% Q13
Per the key, the rural poor are those \textbf{who live with pollution and health hazards} (from chulha smoke, with no electricity).

%% Q14
``Making available \textbf{locally grown} clean and renewable cooking and lighting fuel can improve the quality of rural life.''

%% Q15
The paragraph analyses the problem and compares crops and fuels using evidence (water use, energy output ratio), so it is \textbf{analytical}.

%% Q16
Inattention in class usually signals \textbf{poor communication} by the teacher: unclear, monotonous or irrelevant presentation.

%% Q17
Communication is a network of \textbf{interactions} between people: a dynamic two-way process, not just messages or information.

%% Q18
Noise, culture and attitude can be barriers. \textbf{Job satisfaction} is not a barrier to communication.

%% Q19
A good communicator \textbf{has mastery over the medium of instruction}, i.e.\ command of language and channel.

%% Q20
The most important characteristics of an effective teacher are \textbf{knowledge of content (i)} and \textbf{effective verbal communication (ii)}.

%% Q21
The numbers alternate cube and square: $2^3=8$, $3^2=9$, $4^3=64$, $5^2=25$, $6^3=216$, $7^2=49$. \textbf{125} breaks the pattern (it should be 25).

%% Q22
$2\times2^3=16$, $2\times3^3=54$, $2\times4^3=128$, so next $2\times5^3=\mathbf{250}$.

%% Q23
B(2) $+5$ G(7) $+4$ K(11) $+3$ \textbf{N(14)} $+2$ P(16).

%% Q24
Each letter moves two places forward: C$\rightarrow$E, A$\rightarrow$C, T$\rightarrow$V. SUN $\rightarrow$ \textbf{UWP}.

%% Q25
The differences fall by one: $-8, -7, -6, -5, -4$. So $69-4=\mathbf{65}$.

%% Q26
31 December 2004 is 23 days before 23 January 2005. $23=3\times7+2$, so it falls two days before Sunday: \textbf{Friday}.

%% Q27
George Bush (USA), Gen.\ Musharraf (Pakistan) and APJ Abdul Kalam (India) were then \textbf{heads of state (Presidents)}. \textbf{Tony Blair} was UK \emph{Prime Minister} (head of government), so he is the odd one out.

%% Q28
Shyam: 16--19. Brother: 19--21. The common date is the \textbf{19th}.

%% Q29
1$\rightarrow$B, 2$\rightarrow$C, \ldots, so $n\rightarrow$ the $(n+1)$th letter. 9 $\rightarrow$ 10th letter $=$ \textbf{J}.

%% Q30
A $=$ C and B $=$ C (each individually equals C), and C $=$ D. So A $=$ B $=$ C $=$ D: \textbf{all have the same weight}.

%% Q31 ctx
Answer the questions 31 to 35 from the data given in the following table. In this table the runs made by the five batsmen in five matches are reported.
\begin{center}\small
\begin{tabular}{lccccc|c}\toprule \textbf{Batsman} & M1 & M2 & M3 & M4 & M5 & \textbf{Total}\\\midrule
Mohan & 44 & 53* & 26 & 22* & 38 & 183\\ Anand & 22 & 15 & 42 & 13 & 0 & 92\\ Kishor & 72 & 56 & 48 & 63 & 51 & 290\\ Chetan & 32 & 135 & 0 & 182 & 0 & 349\\ Ramesh & 09 & 11 & 17 & 06 & 10 & 53\\\bottomrule
\multicolumn{7}{l}{\scriptsize * Not out. (Totals column added for the solution.)}
\end{tabular}
\end{center}

%% Q31
Totals: Mohan 183, Anand 92, Kishor 290, \textbf{Chetan 349}, Ramesh 53. Chetan made the highest total.

%% Q32
\textbf{Ramesh} scored low in every match (9, 11, 17, 6, 10), so he is consistently poor.

%% Q33
\textbf{Chetan} scored zero twice (Matches 3 and 5); Anand once.

%% Q34
\textbf{Kishor} scored 48--72 in every match: consistently good. Chetan's runs fluctuate wildly.

%% Q35
Mohan's average (runs $\div$ innings, as per the key) $=183\div5=\mathbf{36.60}$.
\tip{In cricket, the batting average divides by \emph{dismissals} ($183\div3=61$). The key uses the simple arithmetic mean.}

%% Q36
The official key gives \textbf{Server}: the web server is the computer that stores web pages and delivers (serves) their content over the Internet on request.

%% Q36 note
On the user's side, the software that receives and displays web pages is the \emph{browser}; the server is the machine that supplies them. The official key is (C).

%% Q37
MS Word is \textbf{application software}: it performs a user task (word processing).

%% Q38
The first page of a website is the \textbf{home page}.

%% Q39
The internet lets a teacher \textbf{update and enrich his or her knowledge} continuously.

%% Q40
A \textbf{modem} connects a computer to a WAN or the Internet via telephone lines. It is not needed inside a LAN, where printers, CPUs and monitors are used.

%% Q41
India won the Hockey World Cup in \textbf{1975} at \textbf{Kuala Lumpur}, Malaysia, beating Pakistan 2--1 in the final.

%% Q42
\textbf{Demography} is the statistical study of \textbf{human population}: size, structure, births, deaths and migration.

%% Q43
\textbf{Ozone} is useful in the stratosphere (it blocks UV) but harmful at ground level (a pollutant causing respiratory problems).

%% Q44
The Department of Environment was created in 1980 and became the Ministry of Environment and Forests in 1985. Among the options, the key gives \textbf{1982}.

%% Q45
\textbf{Carbon dioxide} is the main greenhouse gas responsible for the enhanced greenhouse effect.

%% Q46
Sardar Patel University, M.~S.~University of Baroda and Bhavnagar University are all \textbf{state universities}, so \textbf{none} is deemed.

%% Q47
The main feature of distance education is \textbf{counselling sessions and study centres}, with learning through self-study material and no compulsory attendance.

%% Q48
In 2006, \textbf{Prof.\ Sukhadeo Thorat} was UGC Chairman (2006--2011). Arun Nigavekar preceded him; Yash Pal was Chairman in 1986--91.

%% Q49
Technical education is governed by the \textbf{AICTE}.

%% Q50
The National Knowledge Commission (2005, chaired by Sam Pitroda) was set up to build excellence in the knowledge system for the 21st century. Per the key, its purpose is to \textbf{find the relevance of knowledge to society}.
